\subsection{Fukaya 范畴}\label{sec:12g}
回顾一下也许有益：Fukaya 范畴中出现 $\Ai$ 结合关系\eref{1.2}，很大程度上来自 Deligne–Mumford–Stasheff 紧化 $\overline\RG^{d+1}$ 的递归结构。更确切地，每棵有两个顶点的稳定 $d$ 叶树——按\sr{fig:9.5}{图~9.5}的记号为 $T(m,\allowbreak n)$——都标记一个余维一边界面
\begin{equation}\label{eq:12.21}
\overline\RG^{T(m,n)}\cong\overline\RG^{d+2-m}\times\overline\RG^{m+1}\subset\partial\overline\RG^{d+1}.
\end{equation}
这些面对应于\eref{1.2}中 $1<m<d$ 的各项。\znote{183}{$\Ai$ 关系的几何来源就是模空间边界的分层：$\overline\RG^{d+1}$ 的每个余维一边界面都对应一种把 $d+1$ 个点分成两组的括号化方式，边界算子在该面上的取值正是相邻括号化之差，求和为零即 (1.2)。这也是为什么 Fukaya 范畴的 $\mu^d$ 必须满足无穷多条关系：紧化本身就携带全部高阶数据。$\Ai$ 代数的标准定义（张量余代数上次数 $2-k$ 的余导子 $D$ 满足 $D^2=0$）正是这套分层结构的代数化身。[来源：nLab “A-infinity algebra”，https:/\allowbreak /\allowbreak ncatlab.org/\allowbreak nlab/\allowbreak show/\allowbreak A-infinity+\allowbreak algebra（2026-11-15 访问）。]}

由\eref{12.8}可见，定义 $\Ai$ 结构需要选择各空间 $\RG^{d+1}$ 的定向。选定后，比较\eref{12.21}中的定向，会对每个 $(m,\allowbreak n)$ 给出一个符号。目标是安排这些符号，使它们与\eref{12.10}产生的 Koszul 符号，以及定义 $\mu^d$ 时引入的其他辅助符号一起，最终给出正确版本的\eref{1.2}。从\eref{9.7}将 $\RG^{d+1}$ 定义为 $\operatorname{Conf}_{d+1}(\partial D)$ 关于 $\operatorname{Aut}(D)$ 的商开始。取 $\operatorname{Conf}_{d+1}(\partial D)$ 中由 $z_0,\allowbreak z_1,\allowbreak z_2$ 固定的构形 $(z_0,\allowbreak \ldots,\allowbreak z_d)$ 组成的子空间。$(z_3,\allowbreak \ldots,\allowbreak z_d)\in(\partial D)^{d-2}$ 的自然边界定向赋予该子空间一个定向。另一方面，它是 $\operatorname{Aut}(D)$ 作用的全局切片，因此可经投影与 $\RG^{d+1}$ 识别。将诱导定向赋予 $\RG^{d+1}$。一个简单但繁琐的
\SourcePage{192}
计算表明，\eref{12.21}中 $\overline\RG^{d+2-m}\times\overline\RG^{m+1}$ 的乘积定向，与 $\partial\overline\RG^{d+1}$ 的边界定向相差 $(-1)^\ddagger$，其中
\begin{equation}\label{eq:12.22}
\ddagger=m(d-n)+m+n.
\end{equation}
这里省略该计算。

假设已如\secref{9j}，完成建立初步版本 $\F(M)^{\mathrm{pr}}$ 所需的全部选择。真正 Fukaya 范畴 $\F(M)$ 定义的头几步显而易见：对象为 $M$ 中的正合 Lagrangian 膜，态射空间为
\[
\operatorname{hom}(L_0^\#,L_1^\#)=CF^*(L_0^\#,L_1^\#),
\]
一阶运算为加上符号的 Floer 边界映射，
\[
\mu^1(a)=(-1)^{|a|}\partial(a).
\]
注意，这使用的是对应于底层 Lagrangian 子流形对 $(L_0,\allowbreak L_1)$ 的 Floer 数据。原则上，也可使 Floer 数据及扰动数据依赖于膜结构。目前这只会徒增复杂性，但后面讨论 Fukaya 范畴的某些变体时将采用这种做法，见\secref{14e}。接着，取膜 $(L_0^\#,\allowbreak \ldots,\allowbreak L_d^\#)$，其中 $d\geq2$，固定 $y_0\in\mathcal C(L_0,\allowbreak L_d)$、$y_k\in\mathcal C(L_{k-1},\allowbreak L_k)$。运用\eref{12.8}及 $\RG^{d+1}$ 的给定定向，可知模空间 $\M^{d+1}(y_0,\allowbreak \ldots,\allowbreak y_d)$ 的维数为 $i(y_0)-i(y_1)-\cdots-i(y_d)+2-d$，且相对于 $o(y_0)\otimes o(y_1)^\vee\otimes\cdots\otimes o(y_d)^\vee$ 定向。现在假设 $(r,\allowbreak u)$ 属于零维模空间。则 $i(y_0)=i(y_1)+\cdots+i(y_d)+2-d$，这正是运算 $\mu^d$ 所需的分次；此外，得到选定同构
\begin{equation}\label{eq:12.23}
c_{r,u}:o(y_d)\otimes\cdots\otimes o(y_1)\longrightarrow o(y_0).
\end{equation}
如\eref{12.18}，将 $\mu^d$ 的 $(y_0,\allowbreak \ldots,\allowbreak y_d)$ 系数定义为相应规范化 $|c_{r,\allowbreak u}|_{\K}$ 之和，但额外乘以符号 $(-1)^\dagger$，其中
\begin{equation}\label{eq:12.24}
\dagger=i(y_1)+2i(y_2)+\cdots+di(y_d).
\end{equation}
立即指出，由于\eref{1.3}的约定，底层上同调范畴确实与此前定义的 $H\F(M)$ 一致。\znote{184}{$\mu^1=(-1)^{|a|}\partial$ 与 $(-1)^\dagger$ 都是纯约定：前者来自第 1 节的符号约定 (1.3)，后者把几何上得到的定向符号调整成 (1.2) 的标准形式。改动这些约定只会把整个 $\Ai$ 结构逐项乘以 $\pm1$，得到同构的 $\Ai$ 结构；因此符号选择不影响范畴，只影响书写。Remark 12.4 把这一点说得更明白：$\RG^{d+1}$ 的整体定向与 $(-1)^\dagger$ 是构造中仅有的两处人为选择，几何本身并不挑出优先的约定。[来源：Seidel, Fukaya categories and Picard–Lefschetz theory, §12g，尤见 (12.24) 与 Remark 12.4。]}
\begin{proposition}\label{st:12.3}
各 $\mu^d$ 构成一个 $\Ai$ 结构。
\end{proposition}
\begin{proof}
如上一小节较简单的例子，证明基于粘合论证，从\eref{9.30}中一对孤立点 $(r_k,\allowbreak u_k)$、$k=1,\allowbreak 2$ 出发。只考虑稳定情形 $1<m<d$，有意略去不稳定情形；纳入后者所需的额外论证紧随\eref{12.19}与\eref{12.20}的证明，留给读者作为练习。
\SourcePage{193}
取同构 $c_{r_k,\allowbreak u_k}$，将第二个的输出插入第一个的第 $n+1$ 个输入以复合。从定向空间看，就是将 $o(\tilde y)$ 与其对偶配对。为简洁，将结果记为 $\tilde c$。由上述一般讨论，$\tilde c$ 诱导一维模空间 $\M^{d+1}(y_0,\allowbreak \ldots,\allowbreak y_d)$ 的一个定向。用足够接近零的参数 $\rho$，将 $(r_k,\allowbreak u_k)$ 粘合成其中的点 $(r,\allowbreak u)$。用 Gromov 紧化\eref{9.28}来说，$(r,\allowbreak u)$ 靠近边界点 $(r_1,\allowbreak u_1,\allowbreak r_2,\allowbreak u_2)$，因而存在指向该边界的典范局部定向。两个定向相差符号 $(-1)^{\S}$。确定该符号之后，由构造可知，对所有点对 $(r_k,\allowbreak u_k)$ 求规范化 $(-1)^{\S}|\tilde c|_{\K}$ 的和为零。要由此得到 $(r_1,\allowbreak u_1,\allowbreak r_2,\allowbreak u_2)$ 对
\[
\mu^{d+1-m}(a_d,\ldots,\mu^m(a_{n+m},\ldots,a_{n+1}),a_n,\ldots,a_1)
\]
的 $(y_0,\allowbreak \ldots,\allowbreak y_d)$ 系数的贡献，还须加入\eref{12.24}中的相应符号：
\[
\begin{aligned}
\dagger_1={}&i(y_1)+2i(y_2)+\cdots+(n+1)i(\tilde y)\\
&+(n+2)i(y_{n+m+1})+\cdots+(d+1-m)i(y_d)\\
={}&i(y_1)+2i(y_2)+\cdots+ni(y_n)\\
&+(n+1)(i(y_{n+1})+\cdots+i(y_{n+m})+2-m)\\
&+(n+2)i(y_{n+m+1})+\cdots+(d+1-m)i(y_d),\\
\dagger_2={}&i(y_{n+1})+2i(y_{n+2})+\cdots+mi(y_{n+m}).
\end{aligned}
\]
剩下要证明
\begin{equation}\label{eq:12.25}
\S+\dagger_1+\dagger_2\equiv\sgnmark_n+(2i(y_1)+\cdots+(d+1)i(y_d))\pmod2,
\end{equation}
右侧第一项给出\eref{1.2}的标准符号，第二项则与 $(m,\allowbreak n)$ 无关。结合前面的论证，便建立结合性方程。

为进行所需计算，考虑\eref{12.10}，取 $\RG_1=\RG^{d+2-m}$、$\RG_2=\RG^{m+1}$，但通过对应于 $T(m,\allowbreak n)$ 的映射\eref{9.13}，将 $\RG=(-1;0)\times\RG_1\times\RG_2$ 换成 $\RG^{d+1}$。相应地，令 $\M_{\mathcal S_1}=\M^{d+2-m}$、$\M_{\mathcal S_2}=\M^{m+1}$，并将 $\M_{\mathcal S}$ 换成 $\M^{d+1}$。图的每个角上都有一维实向量空间，按如下方式定向：
\begin{itemize}
\item 右上角取 $\tilde c$ 与 $\RG^{d+1}$ 的定向，再用水平箭头将之转移为 $\lambda^{\mathrm{top}}T\M^{d+1}$ 的元素。按先前讨论，这正是 $\tilde c$ 诱导的模空间定向。
\item 左下角的向量空间与 $\R$ 典范识别，用水平箭头将其转移为右下角的一个平凡化。由 $c_{r_k,\allowbreak u_k}$ 的定义，这等价于将 $\RG^{d+2-m}$ 和 $\RG^{m+1}$ 的定向与 $c_{r_1,\allowbreak u_1}\otimes c_{r_2,\allowbreak u_2}$ 结合。
\end{itemize}
\SourcePage{194}
接着考察这些定向是否与图中的竖直箭头相容。
\begin{itemize}
\item 取左下角模空间的平凡定向，沿作为粘合映射线性化的向上箭头映射，就得到 $\M^{d+1}$ 指向 Gromov 边界的定向。依定义，它与上述定向相差 $(-1)^{\S}$。
\item 右侧一列两个定向的唯一差别，来自相关族的底空间之间的粘合映射\eref{9.13}，所以等于 $(-1)^\ddagger$。
\end{itemize}
最后，整张图在相差已知 Koszul 符号 $(-1)^\diamond$ 的意义下交换，此时
\[
\diamond=m(d-m-1)+m\sum_{k>n+m}i(y_k).
\]
故 $\S\equiv\diamond+\ddagger\pmod2$。直接计算便得到\eref{12.25}。
\end{proof}
\begin{note}\label{st:12.4}
构造中有两项人为选择：$\RG^{d+1}$ 的定向，更确切地，是边界层应给出定向差\eref{12.22}的要求，它在相差无关紧要的整体符号改变的意义下确定这些定向；以及\eref{12.24}。调整这些项，可以得到目前其他任一种 $\Ai$ 结构的符号约定；几何本身并不挑出优先约定。\stref{命题}{12.3}的另一证明见\bcite{60}第 22 节。
\end{note}

\subsection{平移运算与\texorpdfstring{$\Ai$}{A∞} 结构}\label{sec:12h}
现在证明，逐点由\eref{11.36}定义的正合 Lagrangian 膜平移运算，确实对应正确的代数概念。更明确地，它表示相应的 $\Ai$ 模，该模在\secref{3d}中作为\eref{3.2}的特例定义。这一事实在上同调层次的推论，已在 Floer 上同调的讨论中不加证明地陈述，见\eref{12.12}与\eref{12.13}。

考虑一对正合 Lagrangian 膜 $L_0^\#,\allowbreak L_1^\#$，以及后者的平移 $SL_1^\#$。取 $y\in\mathcal C(L_0,\allowbreak L_1)$。一个有用的记号区分，是用 $\tilde y$ 表示与膜对 $(L_0^\#,\allowbreak SL_1^\#)$ 关联的同一条路径 $y$。这意味着除了 $i(y),\allowbreak o(y)$，还有指标与定向空间 $i(\tilde y),\allowbreak o(\tilde y)$。\stref{引理}{11.21}说明 $i(\tilde y)=i(y)-1$，且有典范同构
\begin{equation}\label{eq:12.26}
\sigma_y:o(y)\longrightarrow o(\tilde y).
\end{equation}
回顾该引理的证明，以指标理论表述，\eref{12.26}来自以下事实：定向算子 $D_{H,\allowbreak \tilde y}$ 是 $D_{H,\allowbreak y}$ 与圆盘上指标为 $n-1$ 的算子 $D_{D,\allowbreak \rho}$ 按此顺序的边界连通和。事实上，最好将 $D_{D,\allowbreak \rho}$ 的定义域限制为在 $\partial D$ 的某个固定点消失的函数空间；这使指标降为 $-1$，并同时去掉粘合公式\eref{11.11}中的额外因子 $\lambda^{\mathrm{top}}(F)$。
\SourcePage{195}
基本断言是：若 $u\in\M_Z^*(y_0,\allowbreak y_1)$ 为孤立 Floer 轨线，$\tilde u$ 为视作 $\M_Z^*(\tilde y_0,\allowbreak \tilde y_1)$ 元素的同一映射，则相应同构\eref{12.17}与\eref{12.26}交换，即
\begin{equation}\label{eq:12.27}
\sigma_{y_0}\circ c_u=c_{\tilde u}\circ\sigma_{y_1}.
\end{equation}
两侧都来自将 $D_{H,\allowbreak \tilde y_0}$ 视为由 $D_{Z,\allowbreak u}$（$u$ 的线性化算子）、$D_{H,\allowbreak y_1}$ 与 $D_{D,\allowbreak \rho}$ 粘合而成。两次出现的次序相同，故\eref{12.27}由\eref{11.9}与\eref{11.11}的结合性得到。使用\eref{12.26}的 $\K$ 规范化，再乘以 $(-1)^{i(y)}$，作识别
\begin{equation}\label{eq:12.28}
CF^*(L_0^\#,SL_1^\#)\cong CF^{*+1}(L_0^\#,L_1^\#).
\end{equation}
依定义，它与 $\partial$ 反交换，因而与 $\mu^1$ 交换。\znote{187}{这一小节把“平移”从几何操作提升为代数结构：$S$ 作用在膜上，(12.28) 给出链层面的同构，而 $(-1)^{i(y)}$ 与 $\star=i(y_1)+\cdots+i(y_{d-1})$ 正是使 $\mu^d$ 与之相容所需的修正。用第 3d 节的语言说，$S$ 给每个对象配一个平移对象，(12.28) 及其与 $\mu^d$ 的相容性正是 $\Ai$ 范畴中平移对象应满足的公理；12e 的 (12.12) 只是它在上同调层的影子。若把 $S$ 看成 $\Ai$ 范畴的自等价，则 (12.28) 就是它到一个平移函子的自然同构在链层面的实现。[来源：Seidel, Fukaya categories and Picard–Lefschetz theory, §3d（$\Ai$ 模的定义，(3.2)）与 §12h。]}

再考虑 $(d+1)$ 元膜组 $(L_0^\#,\allowbreak \ldots,\allowbreak L_d^\#)$，其中 $d\geq2$，以及孤立点 $u\in\M^{d+1}(y_0,\allowbreak \ldots,\allowbreak y_d)$。要将 $L_d^\#$ 平移为 $SL_d^\#$，并如前引入记号 $\tilde y_0,\allowbreak \tilde y_d,\allowbreak \tilde u$。对应于\eref{12.27}的陈述是：定向空间之间的多重线性映射 $c_u,\allowbreak c_{\tilde u}$ 满足
\[
\sigma_{y_0}\circ c_u=(-1)^\star c_{\tilde u}\circ(\sigma_{y_d}\otimes\id\otimes\cdots\otimes\id),
\]
其中 $\star=i(y_1)+\cdots+i(y_{d-1})$。论证与前面相同，只是比较两侧时，在复合次序中必须把 $D_{D,\allowbreak \rho}$ 移过 $D_{H,\allowbreak y_1},\allowbreak \ldots,\allowbreak D_{H,\allowbreak y_{d-1}}$，从而产生 Koszul 型符号。因此 $\mu^d$ 与\eref{12.28}交换，这正是一般 $\Ai$ 范畴中平移对象的定义。为看清这一点，注意除 $\star$ 之外，相关符号量还有：由\eref{12.28}对 $y_d,\allowbreak y_0$ 的定义得到的 $i(y_d)$ 与 $i(y_0)=i(y_1)+\cdots+i(y_d)+2-d$，以及来自\eref{12.24}的 $i(y_1)+\cdots+di(y_d)$ 和 $i(y_1)+\cdots+di(\tilde y_d)$。最终所有符号抵消，得到所需结论。
